% GENERATED FILE. Edit canonical lessons, then run npm run generate:books.
% canonical-source-hash: fnv1a64:e19f342e386dd239
% canonical-lessons: IT-C289-lettore, IT-C289-ascoltatore, IT-C289-spettatore, IT-C289-pubblico, IT-C289-editoriale

\chapter{Chi legge, chi guarda}
\label{ch:it-chi-legge-chi-guarda}
\hlchaptermodality{289}

\begin{chapteropening}
\textbf{By the end of this chapter:} \emph{I can talk about the people who follow the news: readers, listeners, viewers, the public, and the editorial.}

Complete the last lesson of chapter 289: I can talk about the people who follow the news: readers, listeners, viewers, the public, and the editorial.
\end{chapteropening}

\section[il lettore]{il lettore — the reader}

\label{lesson:IT-C289-lettore}



\begin{quote}
Before the new one: say the Italian for to spread, then the Italian for to go into in depth.
\end{quote}

\subsection*{You'll want to know: il lettore}
\textbf{il lettore} — \textquotedblleft{}the reader\textquotedblright{}.

\begin{quote}
Il giornale risponde alle lettere dei lettori.
\end{quote}

\begin{quote}
\emph{The paper answers its readers' letters.}
\end{quote}

A woman is \emph{la lettrice}. Also a player: \emph{il lettore DVD}.

\begin{cousinweb}
\emph{lettore} is inherited from Latin \emph{lēctor} \textquotedblleft{}a reader\textquotedblright{} (EtymDB), from \emph{legere} \textquotedblleft{}to read\textquotedblright{}.
\end{cousinweb}

\begin{grammarlens}[title={What you've built}]
One more word for telling, explaining and describing at B1.
\end{grammarlens}

\subsection*{Your turn}
\begin{itemize}
\raggedright
  \item \emph{Say it:} \emph{il lettore}
  \item \emph{Say it:} \emph{il lettore}, once more
  \item \emph{Say it:} a sentence of your own with \emph{il lettore}
  \item \emph{Recall:} say the Italian for to spread, then the Italian for to go into in depth, then say \emph{il lettore} again
\end{itemize}

\begin{tcolorbox}[breakable,colback=teal!4,colframe=teal!35!black,title={Before you move on}]
What does il lettore mean? (\textquotedblleft{}The reader\textquotedblright{}.) Say it once more.
\end{tcolorbox}
\section[l'ascoltatore]{l'ascoltatore — the listener}

\label{lesson:IT-C289-ascoltatore}



\begin{quote}
Before the new one: say the Italian for to go into in depth, then the Italian for the reader.
\end{quote}

\subsection*{You'll want to know: l'ascoltatore}
\textbf{l'ascoltatore} — \textquotedblleft{}the listener\textquotedblright{}.

\begin{quote}
Gli ascoltatori possono telefonare alla radio.
\end{quote}

\begin{quote}
\emph{Listeners can phone the radio station.}
\end{quote}

A woman is \emph{l'ascoltatrice}.

\begin{cousinweb}
\emph{ascoltatore} is built on \emph{ascoltare} \textquotedblleft{}to listen\textquotedblright{}, from Vulgar Latin \emph{ascultāre}, a form of Latin \emph{auscultāre} (Lewis \& Short).
\end{cousinweb}

\begin{grammarlens}[title={What you've built}]
One more word for telling, explaining and describing at B1.
\end{grammarlens}

\subsection*{Your turn}
\begin{itemize}
\raggedright
  \item \emph{Say it:} \emph{l'ascoltatore}
  \item \emph{Say it:} \emph{l'ascoltatore}, once more
  \item \emph{Say it:} a sentence of your own with \emph{l'ascoltatore}
  \item \emph{Recall:} say the Italian for to go into in depth, then the Italian for the reader, then say \emph{l'ascoltatore} again
\end{itemize}

\begin{tcolorbox}[breakable,colback=teal!4,colframe=teal!35!black,title={Before you move on}]
What does l'ascoltatore mean? (\textquotedblleft{}The listener\textquotedblright{}.) Say it once more.
\end{tcolorbox}
\section[lo spettatore]{lo spettatore — the viewer, the spectator}

\label{lesson:IT-C289-spettatore}



\begin{quote}
Before the new one: say the Italian for the reader, then the Italian for the listener.
\end{quote}

\subsection*{You'll want to know: lo spettatore}
\textbf{lo spettatore} — \textquotedblleft{}the viewer, the spectator\textquotedblright{}.

\begin{quote}
La partita ha avuto dieci milioni di spettatori.
\end{quote}

\begin{quote}
\emph{The match had ten million viewers.}
\end{quote}

A woman is \emph{la spettatrice}.

\begin{cousinweb}
\emph{spettatore} is a learned borrowing from Latin \emph{spectātor} \textquotedblleft{}an onlooker\textquotedblright{} (EtymDB), from \emph{spectāre} \textquotedblleft{}to watch\textquotedblright{}.
\end{cousinweb}

\begin{grammarlens}[title={What you've built}]
One more word for telling, explaining and describing at B1.
\end{grammarlens}

\subsection*{Your turn}
\begin{itemize}
\raggedright
  \item \emph{Say it:} \emph{lo spettatore}
  \item \emph{Say it:} \emph{lo spettatore}, once more
  \item \emph{Say it:} a sentence of your own with \emph{lo spettatore}
  \item \emph{Recall:} say the Italian for the reader, then the Italian for the listener, then say \emph{lo spettatore} again
\end{itemize}

\begin{tcolorbox}[breakable,colback=teal!4,colframe=teal!35!black,title={Before you move on}]
What does lo spettatore mean? (\textquotedblleft{}The viewer, the spectator\textquotedblright{}.) Say it once more.
\end{tcolorbox}
\section[il pubblico]{il pubblico — the audience, the public}

\label{lesson:IT-C289-pubblico}



\begin{quote}
Before the new one: say the Italian for the listener, then the Italian for the viewer.
\end{quote}

\subsection*{You'll want to know: il pubblico}
\textbf{il pubblico} — \textquotedblleft{}the audience, the public\textquotedblright{}.

\begin{quote}
Il pubblico ha applaudito per dieci minuti.
\end{quote}

\begin{quote}
\emph{The audience clapped for ten minutes.}
\end{quote}

As an adjective, public: \emph{un giardino pubblico}.

\begin{cousinweb}
\emph{pubblico} is a learned borrowing from Latin \emph{pūblicus} \textquotedblleft{}of the people, public\textquotedblright{} (Lewis \& Short).
\end{cousinweb}

\begin{grammarlens}[title={What you've built}]
One more word for telling, explaining and describing at B1.
\end{grammarlens}

\subsection*{Your turn}
\begin{itemize}
\raggedright
  \item \emph{Say it:} \emph{il pubblico}
  \item \emph{Say it:} \emph{il pubblico}, once more
  \item \emph{Say it:} a sentence of your own with \emph{il pubblico}
  \item \emph{Recall:} say the Italian for the listener, then the Italian for the viewer, then say \emph{il pubblico} again
\end{itemize}

\begin{tcolorbox}[breakable,colback=teal!4,colframe=teal!35!black,title={Before you move on}]
What does il pubblico mean? (\textquotedblleft{}The audience, the public\textquotedblright{}.) Say it once more.
\end{tcolorbox}
\section[l'editoriale]{l'editoriale — the editorial, the leading article}

\label{lesson:IT-C289-editoriale}



\begin{quote}
Before the new one: say the Italian for the viewer, then the Italian for the audience.
\end{quote}

\subsection*{You'll want to know: l'editoriale}
\textbf{l'editoriale} — \textquotedblleft{}the editorial, the leading article\textquotedblright{}.

\begin{quote}
Il direttore scrive l'editoriale ogni domenica.
\end{quote}

\begin{quote}
\emph{The editor writes the editorial every Sunday.}
\end{quote}

The opinion piece of the paper's editor, \emph{il direttore}.

\begin{cousinweb}
\emph{editoriale} is usually said to be modelled on English \emph{editorial}, built on Latin \emph{ēditor} \textquotedblleft{}one who puts forth\textquotedblright{}, from \emph{ēdere} \textquotedblleft{}to give out\textquotedblright{} (Lewis \& Short).
\end{cousinweb}

\begin{grammarlens}[title={What you've built}]
One more word for telling, explaining and describing at B1.
\end{grammarlens}

\subsection*{Your turn}
\begin{itemize}
\raggedright
  \item \emph{Say it:} \emph{l'editoriale}
  \item \emph{Say it:} \emph{l'editoriale}, once more
  \item \emph{Say it:} a sentence of your own with \emph{l'editoriale}
  \item \emph{Recall:} say the Italian for the viewer, then the Italian for the audience, then say \emph{l'editoriale} again
\end{itemize}

\begin{tcolorbox}[breakable,colback=teal!4,colframe=teal!35!black,title={Before you move on}]
What does l'editoriale mean? (\textquotedblleft{}The editorial, the leading article\textquotedblright{}.) Say it once more.
\end{tcolorbox}
